%! Characteristics of Logarithmic Functions — Unit 8, Lesson 8.0 — Exit Ticket (Answer Key)
\documentclass[10pt]{article}
\usepackage{algebra2-article}
\usepackage{algebra2-key}   % pulls in -boxes; do NOT also load -boxes
\newcommand{\TallMath}[1]{$\displaystyle #1\rule[-1.4em]{0pt}{3.2em}$}
\newcommand{\lgrid}[4]{%
  \draw[step=1, linegray!40, very thin] (#1,#3) grid (#2,#4);
  \draw[->, semithick, charcoal] (#1,0)--(#2,0) node[below right, font=\scriptsize]{$x$};
  \draw[->, semithick, charcoal] (0,#3)--(0,#4) node[left, font=\scriptsize]{$y$};}
% pgfmath has no log_b, so the curve is drawn by parameterizing on y: x = b^(y-k) + h.
\newcommand{\logcurve}[5]{\draw[very thick, forest, domain=#4:#5, samples=120, smooth]
  plot ({pow(#1,\x-(#3))+(#2)},{\x});}

\begin{document}

\pageheader{Unit 8, Lesson 8.0}{Exit Ticket --- Answer Key}
\namedateperiod

\vspace{0.08in}

No notes. The graph below is $f(x) = \log_2(x+4)$. Use it for questions 1 and 2.
\begin{center}
\begin{tikzpicture}[scale=0.44]
  \lgrid{-6}{9}{-3}{4}
  \draw[navy, dashed, semithick] (-4,-3)--(-4,4);
  \begin{scope}
    \clip (-6,-3) rectangle (9,4);
    \logcurve{2}{-4}{0}{-3}{3.70}
  \end{scope}
  \fill[charcoal] (-3,0) circle (3pt) node[below right, font=\scriptsize]{$(-3,0)$};
  \fill[charcoal] (-2,1) circle (3pt) node[above left, font=\scriptsize]{$(-2,1)$};
  \fill[charcoal] (0,2) circle (3pt) node[above left, font=\scriptsize]{$(0,2)$};
  \fill[charcoal] (4,3) circle (3pt) node[above left, font=\scriptsize]{$(4,3)$};
\end{tikzpicture}
\end{center}

\begin{enumerate}[leftmargin=1.8em, itemsep=8pt]
  \item \textbf{Read the features.} Write the domain and range in interval notation.

        Domain: \ans{$(-4,\infty)$} \quad Range: \ans{$(-\infty,\infty)$} \quad
        Vertical asymptote: \ans{$x=-4$}

        \vspace{3pt}
        $x$-intercept: \ans{$(-3,0)$} \quad $y$-intercept: \ans{$(0,2)$} \quad
        $f$ is increasing on \ans{$(-4,\infty)$}

        \vspace{3pt}
        Absolute extrema: \ans{none}. \quad End behavior: as $x\to+\infty$,
        $f(x)\to$ \ans{$+\infty$}; \; as $x\to$ \ans{$-4^{+}$}, $f(x)\to$ \ans{$-\infty$}.

  \item \textbf{Explain.} Why does $f$ have a restricted domain, while $y=2^{x}$ does not?
        And why is the boundary of that domain an \emph{asymptote} rather than an endpoint like the
        radical graphs in Unit 6?
        \ansline{$f$ is a logarithm, so its \emph{argument} must be positive: $x+4>0$ forces
        $x>-4$. There is no exponent that makes $2^{\square}$ zero or negative, so a logarithm has
        nothing to return for a non-positive argument. The exponential $y=2^x$ has no such problem
        --- every real number is a legal exponent, so its domain is all reals.\\[3pt]
        The inequality is \emph{strict}: $x=-4$ makes the argument $0$, which is still illegal. So
        the boundary input is \emph{missing} rather than included, and the graph can only approach
        the line $x=-4$ without ever landing on it. A square root allowed radicand $\ge 0$, so its
        boundary input \emph{was} legal and the graph stood on it as a solid endpoint.}

  \item \textbf{SOL-style multiple choice.} What is the domain of $h(x) = \log_5(3-x)$?
        \begin{enumerate}[label=\Alph*., leftmargin=1.9em, itemsep=1pt]
          \item $x > 3$
          \item $x < 3$ \quad \textcolor{keyred}{\textbf{$\leftarrow$ correct}}
          \item $x \le 3$
          \item all real numbers
        \end{enumerate}
\end{enumerate}

\begin{teachernote}
\textbf{Answer 3:} B. Solve $3-x>0 \Rightarrow -x>-3 \Rightarrow x<3$ (dividing by $-1$ flips the
inequality). Each distractor breaks exactly one idea: \textbf{A} solves the inequality but forgets
the flip; \textbf{C} gets the direction right but uses the Unit~6 radical bracket --- a logarithm
needs a \emph{strict} inequality, because the argument may not equal $0$; \textbf{D} treats a
logarithm as if its argument were unrestricted, the exponential-vs-logarithm confusion.

\textbf{Sort the collected tickets into four piles} to decide how Lesson~8.1 opens:
\emph{got it} / \emph{sign-flip slip} (item~3 answered A) / \emph{asymptote-as-endpoint} (square
brackets in item~1, or item~3 answered C) / \emph{domain confusion} (item~2 cannot say why the
argument must be positive, or item~3 answered D). The third pile is the most common on this lesson
and the cheapest to fix --- one minute on ``$>0$, not $\ge 0$'' clears most of it. The fourth pile
matters most: Lesson~8.1 re-derives the domain algebraically from the definition of a logarithm, so
if that pile is large, reopen Warm-Up item~1(f) before starting the log~$\rightleftarrows$~exponential
conversions.
\end{teachernote}

\end{document}
